\section{The one-cell logarithm and harmonic twists}
\label{subsec:interacting-one-cell}

For a fixed cell $c$ and real twist $\theta$, let
\begin{equation}\label{eq:interacting-free-coarse-log}
 g_D^\theta(t)=\tr\log\left(1+D_*(c_\theta)^{-1}
                   [D_*(c_\theta+t\widehat c_c)-D_*(c_\theta)]\right).
\end{equation}
The logarithm here is the matrix power series.
Set
\begin{align}
 C_{\rm fr}&=2592C_dK_0C_\beta^2
       \left(M^2(D_\beta^{\rm buf})^3L^4h^{-2}+8M^5h\right),
                                                   \notag\\
 B_{\rm wr}&=16\exp(6)C_dM^3K_0p_\beta \mathrm e^{-r_*},
 \qquad B_1=16g_1C_K'p_\beta e_*^{2\varepsilon},       \notag\\
 B_0&=C_{\rm fr}e_*^{4\varepsilon}+B_{\rm wr}+B_1 .
                                \label{eq:interacting-free-log-constants}
\end{align}

\begin{proposition}\label{prop:interacting-one-cell}
Suppose \eqref{eq:interacting-factor-conditions}--%
\eqref{eq:interacting-small-complex-mass} hold,
$r_*>31L+16$, and
\begin{equation}\label{eq:interacting-one-cell-conditions}
 2g_1C_K'e_*^{2\varepsilon}\le1/2,\quad
 B_0\le1/4,\quad
 \tau_s:=c_\tau\ell_*\ge1,\quad
 52B_1\mathrm e^{-\tau_s}\le e_*^{1/4-9\varepsilon}.
\end{equation}
The free ratio $N_0$ has the logarithm
\begin{equation}\label{eq:interacting-one-cell-free-log}
 \ell_0(t)=\sum_X\mathfrak P_*(X,t\widehat c_c)+g_D^0(t),
 \qquad N_0(t)=\mathrm e^{\ell_0(t)} ,
\end{equation}
analytic on $|t|<t_x^+$. It is independent of
$c$ and of the admitted placement, is even
and real on the real axis, and satisfies
\begin{align}
 \left|\sum_X\mathfrak P_*(X,t\widehat c_c)\right|
        &\le C_{\rm fr}e_*^{-2\varepsilon}|t|^2+B_{\rm wr},
                                  \label{eq:interacting-P-probe-gain}\\
 |\ell_0(t)|&\le B_0\qquad(|t|\le t_x).
                                  \label{eq:interacting-free-log-bound}
\end{align}
The one-cell interacting error
$w_c(t)=\mathcal G(te_c;\theta,A)-\ell_0(t)$
is analytic, vanishes at zero, and obeys
\begin{equation}\label{eq:interacting-one-cell-error}
 |w_c(t)|\le3e_*^{1/4-9\varepsilon}
                         \qquad(|t|<t_x^+).
\end{equation}
\end{proposition}

\begin{proof}
Write $K=D_*(c_\theta+t\widehat c_c)-D_*(c_\theta)$.
The real covariance and the integrated insertion
bounds give
\[
 \|D_*(c_\theta)^{-1}K\|_2\le
           2g_1C_Ke_*^{2\varepsilon}\le1/2,\qquad
 \|D_*(c_\theta)^{-1}K\|_1
           \le4g_1C_Kp_\beta e_*^{2\varepsilon}.
\]
Thus the series in \eqref{eq:interacting-free-coarse-log}
converges, is analytic, vanishes at zero, and
has modulus at most $B_1$. Triangularizing
the matrix identifies its exponential with
$\det D_*(c_\theta+t\widehat c_c)/\det D_*(c_\theta)$.
Combine this with
\eqref{eq:interacting-determinant-telescope}.
The lattice symmetries carrying one cell
probe to another preserve the fine determinant,
so the resulting fine ratio is $N_0$.

The bound \eqref{eq:interacting-P-probe-gain}
requires a cancellation before summing the levels.
Put $i=k_*-j$ and $B_j=e^{-1}\widehat{\mathfrak a}_c^{[j]}$.
Rescaling gives
\begin{equation}\label{eq:interacting-probe-derivative}
 e_j|\partial B_j|_\infty\le
                         a_j^2a_*^{-3/2}C_\beta .
\end{equation}
Let $\mathcal Q_j$ be the $M$-cubes containing
endpoints of the supporting bonds of $B_j$.
The buffered support cube yields
\begin{equation}\label{eq:interacting-probe-cube-count}
 |\mathcal Q_j|\le
 4\left((D_\beta^{\rm buf})^3M^{-3}a_j^{-3}+8\right),
 \qquad |\mathcal Q_j|\le L^{3i}p_\beta .
\end{equation}
By locality only polymers containing a cube
of $\mathcal Q_j$ contribute to
\eqref{eq:interacting-P-telescope}.

For such a polymer $X$, put $d=d_M(X)$ and
$s_X=M(d+3)$. Suppose first that $2s_X<\ell_j$.
All bonds meeting $X$ lie in a coordinate
cube of radius $M(d+2)+1$ around a chosen
site $\tilde c$ of $X$. The strict inequality
$2[M(d+2)+1]+\eta_j<2s_X<\ell_j$
makes this cube a lift to the lattice.
Define a real function there by
$\lambda(c)=\sum_\nu B_{j,\nu}(\tilde c)(c-\tilde c)_\nu$,
and extend it by zero outside. On bonds meeting
$X$, its gradient is the constant value
$B_j(\tilde c)$. If $[\,\cdot\,]_X$ restricts
a connection to those bonds, then
\[
 B_X=[B_j-d\lambda]_X,\qquad
 |B_X|_\infty\le3s_X|\partial B_j|_\infty.
\]
The last inequality follows by integrating
the first differences along coordinate paths
inside the lifted cube.

Real gauge invariance must be used before
continuing to the complex gauge parameter.
On the convex open set where
$|[tB_j-zd\lambda]_X|_\infty<e_j^{-1+\varepsilon}$,
the analytic function
$z\mapsto E_j^d(X,[tB_j-zd\lambda]_X)$
equals $E_j^d(X,tB_j)$ for real $z$ near zero:
the unrestricted connection is related to
$tB_j$ by a real gauge transformation, and
locality permits the restriction. The identity
theorem gives this equality throughout that
set. In particular it holds at $z=t$
whenever $|t|<\rho_X$, where
$\rho_X=e_j^{-1+\varepsilon}/(3s_X|\partial B_j|_\infty)$.

The function
$z\mapsto E_j^d(X,zB_X)-E_j^d(X,0)$ is even,
bounded by $2C_dM^3\mathrm e^{-(\kappa+1)d}$ on
$|z|<\rho_X$, and has a double zero at zero.
Apply the maximum principle after dividing
by $z^2$. For $|t|\ge\rho_X$ the original
primitive bound gives the same inequality
with the factor $(|t|/\rho_X)^2\ge1$.
Thus, in both cases,
\begin{equation}\label{eq:interacting-local-constant-removal}
 |E_j^{\rm det}(X,tB_j)|\le
 18C_dM^5C_\beta^2(d+3)^2\mathrm e^{-(\kappa+1)d}
          |t|^2e_j^{-2\varepsilon}a_j^4a_*^{-3}.
\end{equation}
When $\partial B_j=0$, $B_X=0$ and this contribution
vanishes. This proves the cancellation also
at the chamber first step.

Since $(d+3)^2\mathrm e^{-(\kappa+1)d}\le9\mathrm e^{-\kappa_0d}$,
summing \eqref{eq:interacting-local-constant-removal}
over cubes and polymers gives at level $j$
the bound
\[
 648C_dK_0C_\beta^2|t|^2e_j^{-2\varepsilon}a_*^{-3}
           \left(M^2(D_\beta^{\rm buf})^3a_j+8M^5a_j^4\right).
\]
For the remaining polymers $2s_X\ge\ell_j$,
$d\ge r_*L^i/2-3$. The extra exponential in
the primitive bound then gives at level $j$
\[
 2\exp(6)C_dM^3K_0p_\beta L^{3i}\mathrm e^{-r_*L^i}.
\]
Sum over $i\ge1$, using $a_j=L^{-i}a_*$,
$e_j^{-2\varepsilon}=e_*^{-2\varepsilon}L^{i\varepsilon}$,
$\sum_iL^{-i(1-\varepsilon)}\le1$, and
$a_*^{-2}\le L^4h^{-2}$, $a_*\le h$.
For $r_*>31L+16$, the second sum is bounded
by a constant times $\mathrm e^{-r_*}$.
The stated definitions of $C_{\rm fr}$ and
$B_{\rm wr}$ give \eqref{eq:interacting-P-probe-gain}.
Together with $|g_D^0|\le B_1$ this proves
\eqref{eq:interacting-free-log-bound}.

Evenness and reality follow from the corresponding
fine determinant symmetries. Both $\ell_0(t)$
and $\ell_0(-t)$ vanish at zero and have
exponential $N_0(t)$, so they coincide.
For real $t$, the nonzero real ratio $N_0(t)$
is positive by continuity from one at zero;
the normalized logarithm is therefore real.
Cell and placement independence follows in
the same way from the determinant translation
and reflection symmetries.

It remains to control the twist. For
$|\operatorname{Im}\theta_\mu|\le\tau_s$
the connections $c_\theta+st\widehat c_c$
lie in $\mathcal W^+$, and
$\|D_*(c_\theta)^{-1}\|_2\le g_1$.
There is a positive strip margin
$\ell_*(e_*^{2\varepsilon}-e_*^{3\varepsilon})$:
split the imaginary constant connection
into a part bounded by $c_\tau/e_*$ and a
remaining part absorbed into $\mathfrak X_*$.
The latter has vanishing first differences.
The mass condition bounds its total imaginary
part by $\mu_*/24$. On this larger open strip,
\cite[(73)]{TerminalPressure} gives
the same matrix power-series bound $|g_D^\theta|\le B_1$.
The function is holomorphic there and
$2\pi$-periodic in each twist variable.
For real twists, periodicity is conjugation by
the periodic gauge factor
$\exp(-2\pi i c_\mu/\ell_*)$; the identity
theorem extends it to the strip.

We recall the short Fourier estimate to make the
size of this suppression explicit. A bounded
periodic holomorphic function $g$ in a closed
strip of width $\tau\ge1$ has
$|\widehat g(m)|\le B \mathrm e^{-\tau|m|_1}$:
shift each Fourier contour towards
$-i\tau\operatorname{sgn}(m_\mu)$, first using
an interior width and then taking its limit.
The vertical sides cancel by periodicity.
The nonconstant Fourier sum is bounded by
\[
 B\left[\left(\frac{1+\mathrm e^{-\tau}}{1-\mathrm e^{-\tau}}\right)^3-1\right]
                                      \le25B\mathrm e^{-\tau}.
\]
The last inequality follows by expanding the
bracket divided by $\mathrm e^{-\tau}$ as a power
series with positive coefficients and evaluating
at $\mathrm e^{-\tau}=\mathrm e^{-1}$. Thus the difference
between two real values is at most $52B\mathrm e^{-\tau}$.
Apply this to $g_D^\theta(t)$ to obtain
\begin{equation}\label{eq:interacting-twist-suppression}
 |g_D^\theta(t)-g_D^0(t)|\le52B_1\mathrm e^{-\tau_s}
                                  \le e_*^{1/4-9\varepsilon}.
\end{equation}

Finally the normalized logarithm identity and
\eqref{eq:interacting-factor-definition} give
\[
 w_c(t)=g_D^\theta(t)-g_D^0(t)
                +\log f_{\rm int}(te_c)-\log f_{\rm int}(0).
\]
The principal logarithm of $f_{\rm int}$ is
analytic and has modulus at most
$e_*^{1/4-9\varepsilon}$. Both sides vanish
at zero and exponentiate to the same ratio;
uniqueness of the logarithm justifies the
identity. Equation
\eqref{eq:interacting-twist-suppression} proves
\eqref{eq:interacting-one-cell-error}.
\end{proof}

Let $\phi_0(t)=\sum_{i=0}^4N_0^{(i)}(0)t^i/i!$.
Cauchy's estimate applied to
$|N_0(t)-1|\le2B_0$ on $|t|=t_x=4r$ gives
\begin{equation}\label{eq:interacting-free-jet}
 |c_2|\le r^{-2}/16,\qquad |c_4|\le r^{-4}/16,\qquad
 \|\phi_0-1\|_r\le1/6,\qquad \|\phi_0\|_r\le2.
\end{equation}
The odd coefficients vanish. Expanding $\log N_0$
through degree four gives
$\ell_0^{(4)}(0)=c_4-3c_2^2$, the chamber fourth coefficient
of \cite[Theorem~B]{PositiveProbe}.

